\newpage
\pagestyle{plain}

{\selectlanguage{russian}
\begin{center}
    \Large
    \textbf{Аннотация}
\end{center}

\vspace{0.5cm}

Фотоэлектрохимическое (PEC) расщепление воды является одним из наиболее перспективных
методов производства солнечного водорода, однако экспериментальный поиск
высокоэффективных материалов для фотоэлектродов остаётся медленным и ресурсоёмким
процессом.

В данной работе представлена полноценная система машинного обучения, которая предсказывает
три ключевых показателя эффективности PEC-ячеек — плотность фототока (mA/cm\textsuperscript{2}),
КПД преобразования солнечной энергии в водород (STH, \%) и скорость выделения водорода
($\mu$mol\,ч\textsuperscript{-1}\,см\textsuperscript{-2}) — непосредственно по параметрам
материала и условиям эксперимента.

Набор данных из 676 записей был сформирован на основе 529 опубликованных
экспериментальных результатов и дополнен физически обоснованными оценками и
синтетическими строками из литературных источников. Было разработано 20
физически информированных признаков, включая взаимодействие ширины запрещённой
зоны с pH, прокси-оценку перенапряжения, прокси ионной силы и кинетические члены.

Ансамбль из регрессоров XGBoost и Random Forest обучается отдельно для каждой целевой
переменной с оценкой обобщающей способности методом 5-кратной кросс-валидации.
Наилучшие значения R\textsuperscript{2}: для фототока~— 0.454, для STH~— 0.533,
для скорости выделения H\textsubscript{2}~— 0.366.
По сравнению с необогащённым датасетом показатели выросли на 39, 26 и 46
процентных пунктов соответственно.

Система развёрнута в виде REST API (Flask) и одностраничного веб-приложения,
возвращающего предсказания, оценки согласованности моделей и
классификацию производительности (НИЗКАЯ / СРЕДНЯЯ / ВЫСОКАЯ),
калиброванную относительно обучающего распределения.

\vspace{0.5cm}
\noindent\textbf{Ключевые слова:} фотоэлектрохимическое расщепление воды,
машинное обучение, XGBoost, Random Forest, производство водорода, солнечная энергия,
информатика материалов, конструирование признаков, ансамблевое обучение.
}
